% =================================================================
%  style/fonts-geometry.tex
%  Шрифты, язык, геометрия страницы, базовая вёрстка
%  Требует компиляции XeLaTeX (см. первую строку math_analysis.tex)
% =================================================================

\usepackage{fontspec}
\defaultfontfeatures{Ligatures=TeX}

\usepackage{polyglossia}
\setdefaultlanguage{russian}
\setotherlanguage{english}

% -----------------------------------------------------------------
% Шрифты: Computer Modern Unicode.
% В эталоне подключаются локальные .otf из папки ../fonts/.
% Переключатель ниже позволяет собрать конспект и без этой папки.
% -----------------------------------------------------------------
\newif\ifuselocalfonts
% Автовыбор: если шрифты лежат рядом — берём их, иначе системные.
\IfFileExists{../fonts/cmunrm.otf}{\uselocalfontstrue}{\uselocalfontsfalse}
% Принудительно — раскомментируйте нужную строку
% (и поправьте Path=, если папка fonts/ лежит не на уровень выше):
% \uselocalfontstrue
% \uselocalfontsfalse

\ifuselocalfonts
	\setmainfont{cmunrm.otf}[
	Path           = ../fonts/,
	BoldFont       = cmunbx.otf,
	ItalicFont     = cmunti.otf,
	BoldItalicFont = cmunbi.otf
	]
	\setsansfont{cmunss.otf}[
	Path           = ../fonts/,
	BoldFont       = cmunsx.otf,
	ItalicFont     = cmunsi.otf,
	BoldItalicFont = cmunbso.otf
	]
	\setmonofont{cmuntt.otf}[
	Path           = ../fonts/,
	BoldFont       = cmuntb.otf,
	ItalicFont     = cmunit.otf,
	BoldItalicFont = cmuntx.otf
	]
\else
	% Запасной вариант: системные шрифты с кириллицей.
	% Latin Modern сюда не годится — в нём нет кириллических глифов.
	\IfFontExistsTF{CMU Serif}{%
		\setmainfont{CMU Serif}
		\setsansfont{CMU Sans Serif}
		\setmonofont{CMU Typewriter Text}
	}{%
		\setmainfont{DejaVu Serif}
		\setsansfont{DejaVu Sans}
		\setmonofont{DejaVu Sans Mono}
	}
\fi

% -----------------------------------------------------------------
% Геометрия и абзацы
% -----------------------------------------------------------------
\usepackage[left=18mm, top=18mm, right=18mm, bottom=18mm]{geometry}
\usepackage{parskip}
\setlength{\parskip}{6pt}
\setlength{\parindent}{0pt}

\usepackage{enumitem}
\usepackage{tabularx}
\usepackage{array}
\usepackage{ragged2e}
\usepackage{anyfontsize}

% Иллюстрации оформляются как в эталоне по алгоритмам:
%     \begin{figure}[H] \centering ... \caption{...} \label{fig:...} \end{figure}
% [H] (пакет float) ставит рисунок ровно там, где он написан.
\usepackage{graphicx}
\usepackage{float}      % даёт [H] для figure
\usepackage{caption}
\captionsetup{font=small, labelfont=bf, width=0.92\textwidth, skip=6pt}

% -----------------------------------------------------------------
% Интервалы вокруг выключных формул и списков
% -----------------------------------------------------------------
\setlength{\abovedisplayskip}{6pt}
\setlength{\belowdisplayskip}{6pt}
\setlength{\abovedisplayshortskip}{4pt}
\setlength{\belowdisplayshortskip}{4pt}
\setlist[enumerate]{itemsep=3pt, topsep=3pt}

% Рисунки нумеруются внутри \section — согласованно с боксами (Рис. 1.1)
\usepackage{chngcntr}
\counterwithin{figure}{section}
